% Method section.
%
% Same rule as results.tex: every number here is quoted from a generated
% table, and the two tables this section owns -- the protocol and the task
% grid -- are read out of the stored metrics.json files rather than out of the
% configs, because the configs have been edited and the results have not.
% Regenerate with:
%
%     cd cf_worldmodels && python experiments/export_tables.py
%
% Requires booktabs.

\section{The benchmark}
\label{sec:method}

\subsection{What is being measured, and in which component}
\label{sec:method:setup}

A world model factorises into three learned parts: an encoder $E$ that maps an
observation $o_t$ to a latent $z_t$, a transition component $M$ that carries
$(z_t, a_t)$ forward in time, and a decoder $D$ that maps latents back to
pixels. Continual-learning results for model-based agents are almost always
reported at the level of the agent --- returns, or success rates, after a task
switch --- which confounds the three parts with each other and with the policy
trained on top of them. This benchmark measures $M$.

We use the recurrent state-space model of \citet{hafner2019planet} as the common
substrate, in the configuration of \citet{hafner2020dreamer}, with a
convolutional VAE encoder--decoder \citep{ha2018worldmodels}. Observations are
rendered to $64 \times 64$ RGB and scaled to $[0,1]$ in all three families, so
the same architecture spans them. The deterministic state is carried by a
GRU cell,
%
\begin{equation}
  h_t = \mathrm{GRU}\big([z_{t-1}, a_{t-1}], h_{t-1}\big),
  \qquad
  z_t \sim \mathcal{N}\big(\mu_\theta(h_t), \operatorname{diag}
  \sigma^2_\theta(h_t)\big),
  \label{eq:rssm}
\end{equation}
%
and $M = (\mathrm{GRU}, \mu_\theta, \sigma_\theta)$ is the object the forgetting
metrics score. The model is trained on reward-free rollouts from a random
policy, so nothing in the pipeline depends on a task's reward function --- a
property that matters more than it sounds, and Section~\ref{sec:method:families}
returns to it.

A run trains one model on a sequence of tasks $T_1 \ldots T_k$, and measures on
$T_1$ what training on the rest cost. All experiments here use $k = 2$: we write
task~A and task~B, and every metric is an A-versus-A comparison of the model
before and after training on B. Each task contributes held-out episodes that are
never trained on, and the evaluation sets are built from those.

\subsection{Three metrics, and what each one can see}
\label{sec:method:metrics}

Write $M_A$ for the model after task~A and $M_B$ for the same model after task~B.

\paragraph{Prediction fidelity (PF).} The evaluation set $D_A$ is a set of
latent transition triples $(z, a, z')$ drawn from held-out task-A episodes. PF
is the one-step negative log-likelihood the later model assigns to them, minus
the earlier model's:
%
\begin{equation}
  \mathrm{PF} = \mathrm{NLL}(M_B, D_A) - \mathrm{NLL}(M_A, D_A),
  \qquad
  \mathrm{NLL}(M, D) = -\,\mathbb{E}_{(z,a,z') \sim D}
  \big[\log P_\theta(z' \mid z, a)\big].
  \label{eq:pf}
\end{equation}
%
Positive PF means the transition component got worse at task~A. The target is
$z'$, the \emph{next} latent: scoring against $z$ would measure how far the
transition moves the state rather than how well it predicts it.

\paragraph{Rollout divergence (RD).} PF is a one-step quantity, and a
transition model can be locally accurate while compounding error over a horizon.
RD imagines forward from held-out task-A start states under both models, driven
by the same actions, and averages the per-step divergence between their
predicted state distributions:
%
\begin{equation}
  \mathrm{RD} = \frac{1}{H} \sum_{t=1}^{H}
  \mathbb{E}\Big[
  D_{\mathrm{KL}}\big(P_{A}(z_t \mid h^A_{t}) \,\big\|\,
  P_{B}(z_t \mid h^B_{t})\big)\Big].
  \label{eq:rd}
\end{equation}
%
Each model propagates its own predicted mean, so the two rollouts are allowed to
separate; the actions at each step are resampled from the evaluation set's own
action distribution rather than from a Gaussian, which would be off-manifold for
bounded continuous actions and meaningless for one-hot discrete ones. RD is a KL
and is therefore unbounded above, which is a property of the metric that
Section~\ref{sec:results:reporting} has to handle rather than an artefact.

\paragraph{Forward transfer (FT).} Forgetting is only half of a continual
learning result. FT asks what having seen task~A did for learning task~B, against
the counterfactual of a model that trained on B alone under the same budget and
the same data:
%
\begin{equation}
  \mathrm{FT} = \mathcal{L}^{\mathrm{px}}_{B}\big(\text{trained on B from
  scratch}\big) - \mathcal{L}^{\mathrm{px}}_{B}\big(\text{pretrained on
  A}\big),
  \label{eq:ft}
\end{equation}
%
where $\mathcal{L}^{\mathrm{px}}_{B}$ is held-out reconstruction error on task~B
in pixel space. Positive FT means knowledge of A helped. It is measured in
pixels, not in latent NLL, and the reason is not cosmetic: the two models were
trained independently, so their latent spaces are unrelated bases and a latent
likelihood would score one model in coordinates belonging to the other. Pixels
are the one scale the two arms share. One caveat travels with the number:
replay's pretrained arm trains on A and B together during phase~B, so it spends
half its updates on task-A data, and its FT mixes transfer with a halved
effective budget on B.

\paragraph{Declared scope.} $D_A$ is encoded once, by $M_A$, and both models are
scored on those fixed latents. This is what isolates $M$: letting each model
encode with its own drifting VAE would compare every model against its own
moving coordinates rather than against a common evaluation set. The cost is that
PF and RD cannot see drift in the encoder at all, and
Section~\ref{sec:results:encoder} shows that the blind spot is where most of the
forgetting happens. We therefore report, for every run, the held-out task-A
reconstruction error in pixels alongside PF and RD. A second and smaller
restriction: the transitions in $D_A$ are scored from a zero initial hidden
state, so PF measures the one-step map and not the recurrent carry.

\paragraph{On aggregation, and on a metric we withdrew.} An earlier version of
this benchmark combined PF, RD and a policy impact score into a single scalar.
We report PF and RD separately (Section~\ref{sec:results:reporting} gives the
measured reason), and we withdraw the policy impact score: scoring a policy
trained inside the model's imagination requires a controller this benchmark does
not build, and reporting an unimplemented term as a zero is worse than not
announcing it. The suite is PF, RD and FT.

\subsection{Two ways to measure the distance between two environments}
\label{sec:method:distances}

The design varies the distance between task~A and task~B, so that distance has
to be quantified rather than asserted. Two instruments are available and they
are not interchangeable.

The \emph{parametric} distance differences the physics vector
$\phi = [\,\text{gravity}, \text{mass scale}, \text{friction scale}\,]$ directly,
%
\begin{equation}
  d_{\mathrm{param}} = \frac{\lVert \phi_A - \phi_B \rVert_2}
  {\lVert \phi_A \rVert_2},
  \label{eq:dparam}
\end{equation}
%
and is exact where it applies. It applies to four of our nine cells. It is
undefined for a change of layout or of body: the pairs that swap cheetah for
walker leave $\phi$ at its default in both tasks, so
Equation~\ref{eq:dparam} would score the largest construction change in that
family as no change at all.

The \emph{transition} distance is defined wherever transitions can be sampled.
It trains one plain RSSM per environment, from scratch, and takes the expected
divergence between their one-step predictions on shared held-out task-A
transitions:
%
\begin{equation}
  d_{\mathrm{trans}} = \mathbb{E}_{(z,a) \sim D_A}
  \Big[ D_{\mathrm{KL}}\big(P_A(z' \mid z, a) \,\big\|\, P_B(z' \mid z, a)\big)
  \Big].
  \label{eq:dtrans}
\end{equation}
%
It is a property of the task pair and the seed, not of the continual learning
method, so it is computed once per cell and seed and all five methods read the
same value back. The two reference models it needs are the same two that supply
the from-scratch arm of Equation~\ref{eq:ft}, which is why measuring both cost
one pair of extra trainings per cell and seed rather than two: 75 pairs against
the grid's 375 runs, a 20\% surcharge.

Equation~\ref{eq:dtrans} carries a caveat we have not resolved, and it is the
same one that keeps FT in pixel space. The two reference models are trained
independently, so $P_B$ is evaluated on latents produced by $A$'s encoder --- a
basis it has never seen. Part of what
$d_{\mathrm{trans}}$ measures is therefore basis mismatch rather than
dynamics mismatch, and its absolute magnitudes are not comparable across
families for that reason (Table~\ref{tab:tasks} puts MiniGrid near $20$ and
DMControl near $5$). We use it only to rank cells within a family and to rank
the nine cells against each other, and Section~\ref{sec:results:predictors}
reports where that ranking holds and where it fails.

\subsection{Environment families and distance levels}
\label{sec:method:families}

\begin{table}[t]
  \centering
  % The widest table in the paper: seven columns, two of them environment
  % names with their physics. Smaller type and tighter columns rather than
  % \resizebox, which would set this table at a font size nothing else uses.
  \footnotesize
  % 3pt, not 4: at 4pt the seven columns overrun the text block by 2.2pt, and
  % the compiled PDF shows it as booktabs rules poking past the right margin.
  % Each point removed here buys 12pt across the six inter-column gaps.
  \setlength{\tabcolsep}{3pt}
  % Generated by experiments/export_tables.py -- do not edit by hand.
% Regenerate after any change to results/ or to the aggregation.

\begin{tabular}{llllrrr}
\toprule
Family & Level & Task A & Task B & Seeds & $d_{\mathrm{param}}$ & $d_{\mathrm{trans}}$ \\
\midrule
MiniGrid & min & \texttt{Empty-5x5} & \texttt{Empty-8x8} & 10 & -- & 22.35 \\
 & med & \texttt{Empty-8x8} & \texttt{FourRooms} & 10 & -- & 20.64 \\
 & max & \texttt{FourRooms} & \texttt{KeyCorridorS3R1} & 10 & -- & 25.19 \\
\addlinespace
Gymnasium & min & \texttt{HalfCheetah} ($g=9.8$) & \texttt{HalfCheetah} ($g=7$) & 5 & 0.283 & 28.75 \\
 & med & \texttt{HalfCheetah} ($g=9.8$) & \texttt{HalfCheetah} ($g=4$) & 5 & 0.586 & 30.86 \\
 & max & \texttt{HalfCheetah} ($g=9.8$) & \texttt{HalfCheetah} ($g=4$, mass $\times3$, friction $\times0.5$) & 10 & 0.622 & 24.08 \\
\addlinespace
DMControl & min & \texttt{cheetah/run} ($g=9.81$) & \texttt{cheetah/run} ($g=7$) & 5 & 0.284 & 1.87 \\
 & med & \texttt{cheetah/run} & \texttt{walker/run} & 10 & -- & 7.58 \\
 & max & \texttt{cheetah/run} & \texttt{walker/run} ($g=4$, mass $\times3$, friction $\times0.5$) & 10 & -- & 7.46 \\
\bottomrule
\end{tabular}

  \caption{The nine cells, read from the stored results rather than from the
  configuration files. $d_{\mathrm{param}}$ is left empty where differencing a
  physics vector would misdescribe the pair (Section~\ref{sec:method:distances});
  $d_{\mathrm{trans}}$ is the median over the cell's seeds and is comparable
  within a family, not across families.}
  \label{tab:tasks}
\end{table}

Three families cover three regimes an $M$ has to cope with, and
Table~\ref{tab:tasks} lists the pairs. \textbf{MiniGrid}
\citep{chevalierboisvert2023minigrid} is discrete: seven one-hot actions, and
levels built by changing the layout, from a larger empty room to four rooms to a
keyed corridor. \textbf{Gymnasium} is continuous control with variable physics:
six actuators on MuJoCo's HalfCheetah \citep{todorov2012mujoco}, with the three
levels applying progressively stronger perturbations to the same body.
\textbf{DMControl} \citep{tunyasuvunakool2020dmcontrol} is the visually hardest
of the three, and its levels mix a change of body with the same physical
perturbation Gymnasium's maximum level applies.

Gymnasium's three levels are nested by construction, and this turns out to carry
more weight than the rest of the design: both the medium and maximum levels move
gravity from $9.8$ to $4.0$, and the maximum additionally scales mass and
friction. It is strictly the medium perturbation plus more perturbation, applied
to the same task pair. Any claim that the axis orders something monotone can be
checked against those two cells alone, with no appeal to how families compare.

Two constraints shaped the levels and are worth stating because they are easy to
trip over. One world model spans both tasks of a pair, so its GRU has a single
action width and the two tasks must have the same number of actuators --- which
in dm\_control leaves only cheetah and walker. And because the model is trained
on reward-free rollouts, two dm\_control tasks from the same domain are the
\emph{same environment} as far as this benchmark can see: they share physics and
initial-state distribution and differ only in a reward function nothing here
consumes. A level built as \texttt{walker/run} against \texttt{walker/stand}
produces a bit-identical copy of its neighbour, which is how we built one before
we caught it.

\subsection{Methods evaluated}
\label{sec:method:methods}

Five methods, spanning the three standard families of mitigation plus two
bounds.

\begin{itemize}
  \item \textbf{Fine-tuning} --- sequential training with no protection. The
  lower bound.
  \item \textbf{Infinite replay} --- an unbounded buffer of every past task,
  sampled uniformly during phase~B. The approximate upper bound, and the only
  method here that sees task-A data while learning task~B.
  \item \textbf{EWC} \citep{kirkpatrick2017ewc} --- a diagonal-Fisher quadratic
  penalty around the post-task-A weights, $\lambda = 1000$. The Fisher is
  estimated from per-sample gradients of $\log P_\theta(z' \mid z, a)$.
  \item \textbf{Progressive networks} \citep{rusu2016progressive} --- one GRU
  column per task, previous columns frozen, lateral projections from the last
  frozen column into the new one. Like every method here it is evaluated
  without a task oracle: inference runs through the newest column, since no
  task label is available at test time to route by. This differs from the
  original protocol, where the label selects the column --- under oracle
  routing the frozen task-A column would score $\mathrm{PF} = 0$ by
  construction --- and the numbers of
  Section~\ref{sec:results:methods} are statements about the task-agnostic
  reading.
  \item \textbf{UG-MTM} --- our own method: the deterministic transition is
  replaced by a pool of four GRU experts \citep{shazeer2017moe} gated by an
  MC-dropout estimate of epistemic uncertainty \citep{gal2016dropout}. After
  task~A the VAE and the $\mu/\sigma$ projection head are frozen along with the
  previous experts, so only the newly activated expert, the uncertainty head
  and the threshold network train on task~B. It is included as one of five
  methods, not as the paper's result.
\end{itemize}

Two properties of EWC's penalty follow from its definition and are worth naming
before the results rather than after. Its Fisher is defined over
$\log P_\theta(z' \mid z, a)$, which contains no encoder parameter, so it is
\emph{exactly} zero on all twenty VAE tensors: the penalty cannot constrain the
encoder even in principle. It is likewise zero on the GRU's hidden-to-hidden
weights, because the transitions the Fisher is estimated from all start from
$h = 0$. Both zeros are pinned by tests, and
Section~\ref{sec:results:encoder} measures what the first one costs.

\subsection{Protocol}
\label{sec:method:protocol}

\begin{table}[t]
  \centering
  % Generated by experiments/export_tables.py -- do not edit by hand.
% Regenerate after any change to results/ or to the aggregation.

\begin{tabular}{lr}
\toprule
Parameter & Value \\
\midrule
Episodes collected per task (random policy) & 20 \\
Gradient updates per task & 5000 \\
Batch size & 8 \\
Sequence length & 5 \\
Learning rate (Adam) & 1e-03 \\
Latent dimension & 32 \\
GRU hidden dimension & 512 \\
KL weight $\beta$ & 1.0 \\
Held-out task-A episodes & 10 \\
Transitions in $D_A$ & 100 \\
Held-out frames for pixel reconstruction & 200 \\
RD rollout horizon & 15 \\
RD start states & 50 \\
Transitions in EWC's Fisher estimate & 50 \\
EWC $\lambda$ & 1000 \\
UG-MTM MC-dropout samples (training) & 3 \\
Seeds & 5 (0, 1, 2, 3, 4) \\
\bottomrule
\end{tabular}

  \caption{The executed protocol. Generated from the \texttt{protocol} block
  that every result file stores, not transcribed from the configuration.}
  \label{tab:protocol}
\end{table}

Table~\ref{tab:protocol} gives the budget. Per cell and seed: collect episodes
from task~A under a random policy, train, snapshot, collect from task~B, train
again, and evaluate. Rollout collection is shared across the five methods and
the reference pair of that cell and seed, so all six see identical data.

The protocol is declared in the configuration, and every result file records the
block it was produced under. The runner reads all of it and hardcodes none of
it; it skips cells it has already computed, and \emph{refuses} cells produced
under a different protocol or on a different task pair, so two budgets can never
be averaged into one table cell. Table~\ref{tab:protocol} is generated from
those stored blocks rather than from the configuration for the same reason: a
protocol table transcribed by hand describes an intention, and what a benchmark
owes its readers is a description of what ran.

Runs are reproducible from their seed. Seeding covers the global RNGs, cuDNN,
and --- the part that is easy to miss --- each environment's own generator,
including Gymnasium's action space and dm\_control's \texttt{RandomState};
without those the five seeds of a cell train on five different datasets and are
not replicates. Instrumentation is wrapped so that measuring cannot change a
result, which is load-bearing here because UG-MTM's gate keeps MC-dropout active
at evaluation time by design and so consumes the random stream. We reproduced a
cell bit for bit eight days after it first ran.

Each family runs in its own process. This is not tidiness: dm\_control cannot
obtain an OpenGL context in a process where MuJoCo already holds one, and a
30-hour run died at exactly that boundary before the runner was split.
